\section{Discussion}
\label{sec:discussion}

\subsection{Why the abstraction transfers}
The result is consistent with a simple account. A lexical model's features are
drawn from the vocabulary of the projects it trained on; at test time on an
unseen repository, much of its input is out-of-vocabulary, and what remains
correlates with project identity rather than with the weakness. The abstracted
model's alphabet is fixed by the C grammar and the operation taxonomy, so its
input distribution is, by construction, the same on every project. The
stability we observe supports this: fold-to-fold AUC deviation is roughly half
that of the lexical model, which is what one expects if lexical variance is
driven by accidental vocabulary overlap between train and test.

This also explains why the effect is larger on ranking (AUC) than at the top
of the ranking (AUPRC). The abstraction supplies a consistent notion of
\emph{what kind of operation} a node performs, which is enough to order
functions by risk; it discards the specific API, which is often what
distinguishes a genuinely dangerous call from a benign one of the same type.

\subsection{Validating an abstraction is not optional}
The most transferable methodological point concerns RQ1. Program abstractions
are routinely introduced as preprocessing and never evaluated: a taxonomy is
listed, a mapping asserted, and downstream accuracy is reported. Our
experience argues this is unsafe. Two defects in our own mapping --- substring
matching that classified writes and predicates as memory operations, and a
formatter class that absorbed 655 stream-I/O call sites --- were invisible to
every downstream metric and were caught only by annotating call sites and
computing per-class precision. Both would have inflated the very statistic we
use to motivate the taxonomy.

The cost is modest: a few hundred annotated names. We would encourage the
practice as a default for work that introduces a code abstraction, in the same
way that a static analysis is expected to report precision and recall.

\subsection{Negative results}
We implemented a learned edge-weighting module and an affinity-weighted
federated aggregation, and expected both to help. Neither improves on its
ablation. We report them because the alternative --- keeping components a
paper names but its data does not support --- is how unreproducible results
enter the literature, and because the negative result is itself informative:
it locates the effect in the representation rather than in the machinery built
around it.

The taxonomy-refinement result has a concrete explanation worth recording. The
memory-operation class fires on under $1\%$ of AST nodes, so subdividing it
into allocation, release, copy and set cannot shift an aggregate metric
however well-motivated the distinction is semantically. A finer taxonomy would
need to be paired with an evaluation that weights the nodes it refines.

\subsection{Limitations}
\label{sec:limitations}

\textbf{Graph fidelity.} Our graphs are tree-sitter syntax trees with sibling
and def-use approximation edges, not Code Property Graphs: interprocedural
control flow and points-to-precise dependence are absent. Every system we
compare consumes the same graphs, so the comparison is internally fair, but
the absolute ceiling is bounded by this choice.

\textbf{Scale and language.} Experiments use 8{,}000 functions per corpus,
compact models, and C only. The abstraction is specified over a tree-sitter
grammar and is therefore instantiable for other languages, but we have not
demonstrated that, and a taxonomy validated on C should not be assumed to
transfer to, say, Rust's ownership discipline.

\textbf{Sampling and concentration.} Corpus samples are taken in file order
rather than at random, and project mass is concentrated (Chrome and
\texttt{linux} supply about two thirds of BigVul). We hold those projects
train-only and report it, but a differently constructed sample could yield
different absolute numbers.

\textbf{Label noise.} We inherit the labels of the underlying corpora, whose
accuracy is itself contested in the literature that produced them.

\textbf{Effect size.} The AUPRC improvement is small and its interval's lower
bound is at zero. The claim we defend is that the representation matters more
than the architecture choices it is usually compared against --- not that
cross-project detection is solved.

\subsection{A corollary for federated deployment}
Because statistics computed over the abstracted graphs contain no
project-specific vocabulary, class-conditioned summaries can be shared between
organisations without transmitting code, and perturbed to satisfy differential
privacy at low dimensionality. We report this as a corollary rather than a
contribution: in our experiments, ordinary parameter averaging without any
privacy constraint remained the stronger aggregation strategy, and we found no
evidence that the prototype-based alternative beats it absent a formal privacy
requirement.
